% !TeX document-id = {1c0b4298-276a-4038-8e08-c4a1f4846da7}
% !TeX TXS-program:bibliography = txs:///biber
\documentclass[10pt,a4paper]{article}
\usepackage[T1]{fontenc}
\usepackage{graphicx}
\usepackage{mathtools}
\usepackage{amssymb}
\usepackage{amsthm}
\usepackage{thmtools}
\usepackage{xcolor}
\usepackage{nameref}
\usepackage{hyperref}
\usepackage{color}
\usepackage{float}
\usepackage{enumitem}
\usepackage[
backend=biber,
style=authoryear,
natbib=true
]{biblatex}

\usepackage{geometry}
\addbibresource{bibliography.bib}

\title{EAE6029}
\author{Exercícios sobre Variáveis Instrumentais}
\date{}
\begin{document}
	\maketitle
	No que segue, sempre suponha que as variáveis aleatória relevantes estão definidas no mesmo espaço de probabilidade.
 \paragraph{Exercício 1 (Uma equivalência algébrica)} Mostre que o estimador de Wald visto em aula para a estimação de um efeito causal de uma variável escalar $X$ sobre um resultado $Y$, usando como instrumento uma variável escalar $Z$, é \textbf{numericamente idêntico} ao coeficiente estimado de $X$ no estimador de MQO em dois estágios que instrumenta um intercepto e $X$ com base num intercepto e em $Z$.
 
 \vspace{1em}
 		Para os próximos exercícios, suponha sempre que as variáveis aleatórias têm variância estritamente positiva e finita.
 
	\paragraph{Exercício 2 (Erro de medida)} Considere um modelo causal linear da forma:

$$Y = X^* \beta + U\, ,$$
onde $X^*$ é uma causa escalar de interesse, $U$ são as demais causas de um fenômeno, e $\beta$ é um efeito causal de interesse. 

Suponha que o pesquisador não observe $X^*$, mas somente uma medida ruidosa $X = X^* + V$ de $X^*$, onde o erro de medida $V$ é tal que $\operatorname{cov}(V,X^*) = 0$ e $\operatorname{cov}(V, U)=0$. Suponha, ademais que $\operatorname{cov}(X^*,U)$.

\begin{itemize}
	\item[a] Considere o estimador de MQO de $Y$ num intercepto e $X$, baseado numa amostra aleatória de $(Y,X)$. Mostre que o estimador de MQO do coeficiente associado a $X$ \textbf{não} é consistente para $\beta$. Encontre o limite em probabilidade desse estimador. Qual é sua relação com $\beta$?
	\item[b] Suponha que você também observe uma segunda medida ruidosa de $X^*$, $Z = X^* + \xi$, onde $\operatorname{cov}(\xi, X^*) = 0$, $\operatorname{cov}(\xi, U) = 0$ e $\operatorname{cov}(\xi, V) = 0$. Sugira um estimador consistente para $\beta$. Justifique.	Se o seu interesse reside em fazer previsões de $Y$ com base nas medidas observáveis de $X^*$ para novas observações fora de sua amostra, faz sentido usar esse estimador? Justifique.
\end{itemize}

	\paragraph{Exercício 3 (Instrumentos fracos)} Considere um modelo causal linear da forma:

$$Y = X \beta + U\, ,$$
onde $X$ é uma causa escalar de interesse, $U$ são as demais causas de um fenômeno, e $\beta$ é um efeito causal de interesse. Su

\begin{itemize}
	\item[a] Calcule o limite de probabilidade do estimador de MQO baseado em uma amostra aleatória de $(X,Y)$.
	\item[b] Suponha, agora, que tenhamos acesso a um instrumento $Z$ relevante, mas não necessariamente válido. Calcule o limite de probabilidade do estimador de Wald baseado em uma amostra aleatória de $(X,Y,Z)$. Sob quais condições sobre $ \operatorname{cor}(X,U)$, $\operatorname{cor}(X,Z)$ e $\operatorname{cor}(Z,U)$, temos que $|\operatorname{plim}\hat \beta_{IV} -\beta|  > |\operatorname{plim}\hat \beta_{MQO} -\beta| $? Discuta seus resultados.
	
	\item[c] Suponha agora, que a relação de primeiro estágio entre $X$ e $Z$ é dada, numa amostra aleatória de tamanho $n$, por um modelo preditivo linear da forma:
	
	$$X = \gamma_n Z + V\, , \quad \operatorname{cov}(Z,V)=0$$
	onde $\gamma_n = C/\sqrt{n}$ para algum $C \neq 0$. A ideia é capturar uma situação em que a relevância do instrumento é da ordem da incerteza amostral.
	
	Suponha que $Z$ é válido, i.e. $\operatorname{cov}(Z,U)=0$. Mostre que:
	
	$$\hat{\beta}_{IV} \overset{d}{\to} \beta + \frac{N_1}{C\mathbb{V}[Z] + N_2}\, , $$
	onde $$\footnotesize \begin{pmatrix}
	N_1 \\
	N_2
	\end{pmatrix}  \sim \mathcal{N}\left(\begin{bmatrix}
	0 \\
	0
	\end{bmatrix}, \begin{bmatrix}
	\mathbb{E}[(Z-\mathbb{E}[Z])^2 (U  - \mathbb{E}[U])^2] & 	\mathbb{E}[(Z-\mathbb{E}[Z])^2 (U  - \mathbb{E}[U])(V-\mathbb{E}[V])] \\
	\mathbb{E}[(Z-\mathbb{E}[V])^2 (U  - \mathbb{E}[U]) (V-\mathbb{E}[V])] & 	\mathbb{E}[(Z-\mathbb{E}[Z])^2 (V  - \mathbb{E}[V])^2] 
	\end{bmatrix}\right)\, ,$$
	
	Conclua que o estimador de variáveis instrumentais, nesse caso, é inconsistente. Mostre que, nesse mesmo regime assintótico, temos que:
	
	$$\hat \beta_{MQO}\overset{p}{\to} \beta + \theta\, ,$$
	onde $\theta = \frac{\operatorname{cov}(U,V)}{\mathbb{V}[V]}$.
	
	\vspace{0.5em}
	Para os próximos dois itens, vamos considerar o regime assintótico do item (c), e, ademais, supor as condições de homocedasticidade:
	
$$\mathbb{E}[(Z-\mathbb{E}[Z])^2 (U  - \mathbb{E}[U])(V-\mathbb{E}[V])]  = \mathbb{V}[Z]\operatorname{cov}(U,Z)\, .$$

$$\mathbb{E}[(Z-\mathbb{E}[Z])^2 (V-\mathbb{E}[V])^2] = \mathbb{V}[Z]\mathbb{V}[V]$$
		
	\item[d] 	Seja $f:\mathbb{R}^2 \mapsto \mathbb{R}$ a função dada por $f(x) = \frac{x_1}{x_2}$. Para uma sequência de vetores aleatórios $\boldsymbol{S}_n$  com valores em $\mathbb{R}^2$ em que a média da variável aleatória na segunda posição é diferente de zero e em que $\lVert \boldsymbol{S}_n - \mathbb{E}[\boldsymbol{S}_n] \rVert = O_p(n^{-1/2})$, use a expansão de Taylor de segunda ordem:
	
	$$f(\boldsymbol{S}_N) - f(\mathbb{E}[\boldsymbol{S}_N]) \approx \nabla_x f(\mathbb{E}[\boldsymbol{S}_N])(S_n -\mathbb{E}[\boldsymbol{S}_n]) + \frac{1}{2} (S_n -\mathbb{E}[\boldsymbol{S}_n])' H(f;\mathbb{E}[\boldsymbol{S_n}]) (S_n -\mathbb{E}[\boldsymbol{S}_n]) + o_p(n^{-1})\, ,$$
	e a aproximação resultante para a esperança:\footnote{Essa aproximação é conhecida na literatura como expansão de Nagar do viés de um estimador. Ela implicitamente supõe que o erro $R_n = o_p(n^{-1})$ da expansão de Taylor acima é tal que $\mathbb{E}[R_n] = o(n^{-1})$. Condições suficientes para que isso seja verdade estão fora do escopo desse curso. }
	
	$$\mathbb{E}[f(\boldsymbol{S}_n) - f(\mathbb{E}[\boldsymbol{S}_n])] =  + \frac{1}{2} \mathbb{E}[(S_n -\mathbb{E}[\boldsymbol{S}_n])' H(f;\mathbb{E}[\boldsymbol{S}_n]) (S_n -\mathbb{E}[\boldsymbol{S}_n]) ] + o(n^{-1})\, ,$$
	para aproximar, no regime assintótico de (c), o viés do estimador de variável instrumental por:
	
	$$\mathbb{E}[\hat{\beta}_{IV}] - \beta = -\theta \frac{\mathbb{V}[V]}{C^2\mathbb{V}[Z]} + o(n^{-1})$$
	
			\item[e] Considere a estatística $t$ de primeiro estágio, calculada sob homocedasticidade, i.e.
	
	$$\hat{t} = \frac{\sum_{i=1}^n (X_i -\bar{X})(Z_i - \bar{Z})}{\hat{\sigma}_V \sqrt{\sum_{i=1}^n(Z_i-\bar{Z})^2}}\, .$$
	
	Mostre que:
	
	$$\hat t \overset{d}{\to} \mathcal{N}\left(\frac{C \operatorname{sd}(V)}{ \operatorname{sd}(Z)}, 1\right)\, .$$
	
	Calcule a distribuição assintótica da estatística $\hat F = \hat t^2$.

	
	

   \item[f] Com base na estatística $\hat F$, proponha um teste assintoticamente correto, ao nível de significância de 5\%, da hipótese nula de que o viés assintótico do estimador de IV é no máximo 10\% do ``viés assintótico'' do estimador de MQO, i.e., queremos testar
   
   $$H_0: \frac{|\lim_{n \to \infty}\mathbb{E}[\hat{\beta}_{IV}] - \beta |}{|\operatorname{plim}_n \hat{\beta}_{MQO} -\beta|}\leq 0.1\, ,$$
   contra a alternativa de que o viés relativo é maior que 10\%.

\end{itemize}

\paragraph{Exercício 4 (Mercado de Peixe Fulton)} Considere o seguinte modelo causal linear para descrever as curvas de oferta e demanda no período $t$ em um mercado:
$$\log(D_t(p)) = \psi \log(p) + \epsilon^D_t\, ,$$
$$\log(S_t(p)) = \gamma \log(p) + \delta' Z_t + \epsilon^S_t\, , $$
onde $\psi$ é a elasticidade-preço da demanda, e $\epsilon^D_t$ descrevem \textit{shifters} não observáveis da curva de demanda no mercado no período $t$. De forma análoga, $\gamma$ denota a elasticidade-preço da curva de demanda, $\epsilon^S_t$ denotam \textit{shifters} não observáveis da curvam oferta, e $Z_t$ são shifters observáveis da curva de oferta que sabemos não exibirem associação sistemática com os \textit{shifters} da demanda, i.e. $\operatorname{cov}(Z_t, \epsilon^S_t)=0$.

\begin{itemize}
	\item[a] Suponha que os preços e quantidades vendidas $(P_t,Q_t)$ em cada instante de tempo $t$ seja determinado pelo cruzamento das curvas de oferta e demanda, i.e.
	$$Q_t = S_t(P_t)= D_t(P_t)\, .$$
	
	Mostre que, se for este o caso, o melhor preditor linear de $\log(Q_t)$ num intercepto e em $\log(P_t)$ não identificará a elasticidade $\psi$, e que a inclusão de $Z_t$ na predição linear também não garante a identificação.
	
	\item[b] Proponha uma estratégia de variável instrumental para identificar $\psi$ e demonstre a consistência do estimador resultante, baseado em uma amostra aleatória dos mercados ao longo de $T$ períodos.
	
	\item[c] Com base nos dados do mercado de peixe Fulton \citep{Graddy2006}, disponíveis \href{https://users.ox.ac.uk/~nuff0078/Data/fulton.zip}{neste link}, proponha uma estratégia de variável instrumental para identificar a elasticidade da demanda nesse contexto. Implemente a estimação dessa estratégia. Com base em um nível de significância de 5\%, o instrumento é relevante? Você rejeita a hipótese nula de que a demanda tem elasticidade $-1$? \textit{Observação:} no \texttt{R}, use o pacote \texttt{fixest} para implementar a estratégia com erros padrão robustos a heterocedasticidade.
\end{itemize}


\paragraph{Exercício 5 (\textit{Eminent Domain})} Neste exercício, vamos estudar o impacto do \emph{eminent domain} sobre os preços de imóveis nos Estados Unidos\footnote{Este exercício é baseado em \citet{chen2015growth}.}. \emph{Eminent domain} (ou leis de desapropriação) é o mecanismo legal pelo qual o governo pode adquirir compulsoriamente uma propriedade privada. Esse tipo de ação é bastante controverso, e os proprietários podem recorrer à Justiça para tentar reverter a decisão. Quando tribunais federais decidem que uma desapropriação foi ilegal (\emph{decisões favoráveis ao autor}, \emph{pro-plaintiff}), os direitos de propriedade são reforçados e cria-se um precedente que torna mais difícil a utilização futura do \emph{eminent domain}. Nosso objetivo será analisar o efeito causal dessas decisões favoráveis ao autor, proferidas nos tribunais federais de apelação, sobre os preços de imóveis.\footnote{``A foundational understanding of the U.S. federal courts is important to understanding our identification strategy, which relies on the law-making function of common law courts. This making of law occurs because a judge's decisions in current cases become precedent for use in future cases in the same court and in lower courts of the same jurisdiction. There are three layers of federal courts: District, Circuit, and the U.S. Supreme Court. The 94 U.S. District Courts serve as the general trial courts, where a jury is drawn to decide \emph{issues of facts}. If a party appeals the decision, the case goes up to a Circuit Court, which decides \emph{issues of law}; they take facts as given from District Courts and have no juries. The 12 U.S. Circuit Courts, also known as Courts of Appeals or federal appellate courts, are only to hear cases presenting new legal issues (only 10-20\% of District Court opinions are appealed). Cases that reach the Circuit Courts are the more challenging and controversial cases with the greatest likelihood to set new precedent.'' \citep[p. 11-12]{chen2015growth}.} 

\begin{itemize}

\item[a] Importe \texttt{temp\_CSHomePrice.mat.csv}, disponível na subpasta \texttt{dados}. Esse conjunto de dados contém informações sobre preços de imóveis (índice Case-Shiller) e decisões favoráveis ao autor por circuito judicial e período. Observe que, com exceção do índice Case-Shiller em log, todas as variáveis estão em frequência anual. Agregue os dados por ano, calculando a média \emph{dentro de cada circuito-ano} do log do índice de preços.

\item[b] Estime uma especificação linear por MQO do log do índice de preços em um intercepto e  no número total de decisões de desapropriação em apelação favoráveis ao autor (\texttt{numpro\_casecat\_12}). O que você encontra? Esse coeficiente pode ser interpretado como uma estimativa de um efeito causal? Por quê?

\vspace{0.5em}

Para lidar com possíveis problemas de endogeneidade, \citet{chen2015growth} exploram o fato de que, nos tribunais federais de apelação, os juízes são \emph{designados aleatoriamente} para compor painéis de três juízes que decidem os casos. Cada caso recebe três juízes sorteados de um conjunto (\emph{pool}) de magistrados do circuito, que varia aproximadamente entre 8 e 40 dependendo do tamanho do tribunal. Assim, \emph{condicionalmente às características desse conjunto de juízes}, as características observadas dos painéis sorteados devem satisfazer a restrição de exclusão, pois só podem afetar os preços dos imóveis por meio das decisões judiciais.

\item[c] Faça o \emph{merge} do conjunto de dados \texttt{probs.txt} com sua base. Esse arquivo contém a distribuição das características dos juízes em cada circuito entre 1945 e 2007.

\item[d] Considere o seguinte modelo linear causal:

\begin{equation}
	\ln(\text{Case\_Shiller})_{ct} = \beta_0 + \beta_1 \text{numpro\_casecat\_12}_{ct} + \gamma' W_{ct} + \epsilon_{ct}\, ,
\end{equation}
onde $W_{ct}$ inclui o conjunto completo de controles para as características do pool de juízes\footnotemark; uma variável indicadora para os casos em que não houve processos naquele circuito-ano (\texttt{missing\_cy\_12}); e o número total de decisões de apelação envolvendo desapropriações (\texttt{numcasecat\_12}). Como instrumento, utilizaremos o número de painéis com \emph{pelo menos um juiz Democrata} (\texttt{numpanels1x\_dem})\footnote{A justificativa \emph{a priori} dos autores para essa escolha baseia-se em evidências de que juízes Democratas tendem a ser mais favoráveis aos proprietários em casos de \emph{eminent domain}.}. 
\footnotetext{A lista completa de controles para características do pool judicial é: prob\_1x\_dem, prob\_2x\_dem, prob\_3x\_dem, 
	prob\_1x\_female, prob\_2x\_female, prob\_3x\_female, 
	prob\_1x\_nonwhite, prob\_2x\_nonwhite, prob\_3x\_nonwhite, 
	prob\_1x\_black, prob\_2x\_black, prob\_3x\_black, 
	prob\_1x\_jewish, prob\_2x\_jewish, prob\_3x\_jewish, 
	prob\_1x\_catholic, prob\_2x\_catholic, prob\_3x\_catholic, 
	prob\_1x\_noreligion, prob\_2x\_noreligion, prob\_3x\_noreligion, 
	prob\_1x\_instate\_ba, prob\_2x\_instate\_ba, prob\_3x\_instate\_ba, 
	prob\_1x\_ba\_public, prob\_2x\_ba\_public, prob\_3x\_ba\_public, 
	prob\_1x\_jd\_public, prob\_2x\_jd\_public, prob\_3x\_jd\_public, 
	prob\_1x\_elev, prob\_2x\_elev, prob\_3x\_elev, 
	prob\_1x\_female\_black, prob\_2x\_female\_black, prob\_3x\_female\_black, 
	prob\_1x\_female\_noreligion, prob\_2x\_female\_noreligion, prob\_3x\_female\_noreligion, 
	prob\_1x\_black\_noreligion, prob\_2x\_black\_noreligion, prob\_3x\_black\_noreligion, 
	prob\_1x\_female\_jd\_public, prob\_2x\_female\_jd\_public, prob\_3x\_female\_jd\_public, 
	prob\_1x\_black\_jd\_public, prob\_2x\_black\_jd\_public, prob\_3x\_black\_jd\_public, 
	prob\_1x\_noreligion\_jd\_public, prob\_2x\_noreligion\_jd\_public, prob\_3x\_noreligion\_jd\_public, 
	prob\_1x\_mainline, prob\_2x\_mainline, prob\_3x\_mainline, 
	prob\_1x\_evangelical, prob\_2x\_evangelical, prob\_3x\_evangelical, 
	prob\_1x\_protestant, prob\_2x\_protestant, prob\_3x\_protestant.}

Estime o modelo acima utilizando o instrumento proposto. O instrumento satisfaz a condição de \emph{relevância}? Comente os resultados obtidos.

\item[e] Você considera que o coeficiente estimado em (d) fornece uma estimativa \emph{crível do efeito causal} que queremos analisar? Por quê? Caso considere que não, seria possível utilizar os dados disponíveis para estimar uma \emph{especificação alternativa} que produza estimativas mais confiáveis? O que você encontrou?
\end{itemize}

\paragraph{Exercício 6 (Interpretação Causal do Estimando de 2SLS sob Heterogeneidade dos efeitos de tratamento e controles)} Considere uma tupla de variáveis aleatórias $(Y,D,Z,X)$, definidas num espaço de probabilidade  $(\Omega, \Sigma,\mathbb{P})$, em que a distribuição conjunta de $(Y,D,Z,X)$ reflete a distribuição em uma população de um resultado escalar ($Y$) de interesse, uma causa binária ($D$) de $Y$ observável de interesse, um ``instrumento'' binário $Z$ e um conjunto de causas observáveis $X$ de $Y$ sobre as quais não temos interesse direto, e que tomam valores num conjunto finito $\{x_1,\ldots, x_L\}$. Nós definimos os resultados potenciais $\{Y(d,z)\}_{(d,z)\in \{0,1\}^2}$ correspondentes ao experimento hipótetico de se amostrar um elemento da população de interesse e fixar externamente as causa $(D,Z)$ em $(d,z)$, mantidas as demais causas de $Y$, observáveis ou não, em seus valores originais. Também definimos resultados potenciais $(D(0),D(1))$ correspondentes ao experimento hipotético de se fixar externamente o valor de $Z$ em $0$ ou $1$, mantidas as demais causas de $D$, observáveis ou não, fixas.
Com base nessa definição do modelo causal, segue que: 

$$Y = Y(D,Z) = DZ Y(1,1) +(1-D)ZY(0,1) + D(1-Z)Y(1,0) +(1-D)(1-Z)Y(0,0)$$
$$D = ZD(1) + (1-Z)D(0)\, .$$

Ademais, nós supomos as seguintes restrições:

\begin{enumerate}
	\item (Exclusão) $Y(d,0) = Y(d,1) =: Y(d)$ para $d \in \{0,1\}$.
	\item (Aleatorização) Condiconalmente a $X$, $\{Y(1),Y(0),D(1),D(0)\}$ é independente de $Z$.
	\item (Monotonicidade condicional)  $\mathbb{P}[1 \in \{\mathbb{P}[D(1)\geq D(0)|X] ,\mathbb{P}[D(1)\leq D(0)|X]\}]=1$.
	\item (Primeiro estágio condicional)  $\mathbb{E}[(\mathbb{E}[D|Z,X]-\mathbb{E}[D|X])^2]>0$.
\end{enumerate}

\vspace{0.5em}

No que segue, responda às seguintes perguntas:

\begin{enumerate}
	\item[a] Proveja uma interpretação para cada uma das quatro hipóteses.
	\item[b] Seja $\boldsymbol{p}(X) = \begin{bmatrix}
		\mathbf{1}_{\{x_1\}}(X) \\
		\mathbf{1}_{\{x_2\}}(X) \\
		\vdots \\
		\mathbf{1}_{\{x_L\}}(X) 
	\end{bmatrix}$, um vetor com ``dummies'' indicadoras dos pontos de suporte de $X$. Mostre que o estimador de MQO em dois estágios de $\beta$ no sistema linear:
	
	$$Y = \beta D + \gamma' \boldsymbol{p}(X) + U\, , $$
		$$D = \phi_1'Z \boldsymbol{p}(X)  + \phi_0' \boldsymbol{p}(X) + V\, ,$$
		construído com base numa amostra aleatória de $(Y,D,Z,X)$ é consistente para o parâmetro:
		
		$$\beta^*_{WLATE} = \mathbb{E}[w(X)\tau_{LATE}(X)]$$
		onde, $\tau_{LATE}(X)$ é o LATE condicional, definido como:
		
		$$\tau_{LATE}(X) = \frac{\mathbb{E}[(Y(1)-Y(0))\mathbf{1}_{\{D(1)\neq D(0)\}}|X]}{\mathbb{P}[D(1)\neq D(0)|X]}$$
		e $w(X)$ são pesos tais que $w(X) \geq 0$ quase certamente, e $\mathbb{E}[w(X)]$. Para quais subpopulações definidas pelos valores de $X$ o estimando atribui maior peso ao efeito causal? \textit{Dica:} use o teorema de Frisch-Waugh Lovell para mostrar que o limite de probabilidade do estimador de 2SLS é:
		$$\operatorname{plim}_{n \to \infty} \hat \beta_{n,2SLS} = \frac{\mathbb{E}[(\mathbb{E}[D|Z,X]-\mathbb{E}[D|X])(\mathbb{E}[Y|Z,X]-\mathbb{E}[Y|X])]}{\mathbb{E}[(\mathbb{E}[D|Z,X]-\mathbb{E}[D|X])^2]}$$
	\end{enumerate}
		\vspace{0.5em}
		
Você se incomodou com a estimação do sistema linear dado acima, no sentido de que ele envolve um número grandes de instrumentos efetivos ($L$ variáveis excluídas). Você propõe, então, estimar um sistema mais simples, dado por:

	$$Y = \alpha D + \pi' \boldsymbol{p}(X) + \epsilon\, , $$
$$D = \tau Z   + \psi' \boldsymbol{p}(X) + \xi \, .$$

Um colega, entretanto, tendo lido um artigo recentemente \citep{blandhol2022tsls}, nota que, nesse caso, para obter uma interpretação causal sob heterogeneidade irrestrita dos efeitos de tratamento, há que se modificar 3 e 4, para.

\begin{itemize}
	\item[3'] (Monotonicidade incondicional) $\mathbb{P}[D(1)\geq D(0)] =1$ ou $\mathbb{P}[D(1)\leq D(0)] =1$ .
	\item[4'] (Primeiro estágio') $\mathbb{E}[(Z-\mathbb{E}[Z|X])(D-\mathbb{E}[D|X])] \neq 0$.
\end{itemize}

\begin{itemize}
	\item[c] Interprete as novas condições. Mostre que $3' \implies 3$, mas que a recíproca não é verdadeira.
	\item[d] Encontre a interpretação causal do estimando de 2SLS do novo sistema em termos de uma média poderadas de LATEs condicionais. Como os pesos diferem da representação do estimando em (b)?
	\item[\textbf{Bônus}] Você consegue encontrar uma intuição para a necessidade da condição $3'$ mais forte de heterogeneidade com o primeiro estágio simplificado?
\end{itemize}


\paragraph{Exercício 7 (Heterogeneidade e tratamentos ordenados)} Considere uma tripla de variáveis aleatórias $(Y,D,Z)$, definidas num espaço de probabilidade  $(\Omega, \Sigma,\mathbb{P})$, em que a distribuição conjunta de $(Y,D,Z)$ reflete a distribuição em uma população de um resultado escalar ($Y$) de interesse, uma causa  $D$ de $Y$ observável de interesse, e um ``instrumento'' binário $Z$. Nós supomos que a causa $D$ toma valores no subconjunto finito de números reais $\mathcal{D}:=\{d_1,\ldots, d_K\}$, onde $d_1 < d_2 < \ldots < d_K$. Definimos os resultados potenciais $\{Y(d,z)\}_{(d,z)\in \mathcal{D} \times \{0,1\} }$ correspondentes ao experimento hipótetico de se amostrar um elemento da população de interesse e fixar externamente as causa $(D,Z)$ em $(d,z)$, mantidas as demais causas de $Y$, observáveis ou não, em seus valores originais. Também definimos resultados potenciais $(D(0),D(1))$ correspondentes ao experimento hipotético de se fixar externamente o valor de $Z$ em $0$ ou $1$, mantidas as demais causas de $D$, observáveis ou não, fixas.
Com base nessa definição do modelo causal, segue que: 

$$Y = Y(D,Z) = \sum_{i=1}^K\mathbf{1}_{\{d_i\}}(D)[Z Y(d,1) + (1-Z)Y(d,0)]\, ,$$
$$D = ZD(1) + (1-Z)D(0)\, .$$

Ademais, nós supomos as seguintes restrições:

\begin{enumerate}
	\item (Exclusão) $Y(d,0) = Y(d,1) =: Y(d)$ para todo $d \in \mathcal{D}$.
	\item (Aleatorização) $\{Y(1),Y(0),D(1),D(0)\}$ é independente de $Z$.
	\item (Monotonicidade)  $\mathbb{P}[D(1)\geq D(0)]=1$.
	\item (Primeiro estágio)  $\mathbb{E}[(\mathbb{E}[D|Z]-\mathbb{E}[D])^2]>0$.
\end{enumerate}

Mostre que, sob as hipóteses (1) a (4), temos que o estimando de Wald do efeito causal de $D$ sobre $Y$, instrumentado por $Z$, identifica:

$$\beta^*_{\text{Wald}} = \sum_{i=1}^{K-1} \omega_{i}\beta_{\text{LATE}}(d_i,d_{i+1})\, , $$
onde os pesos $\omega_i$ são não negativos e somam um, e $$\beta_{\text{LATE}}(d_i,d_{i+1}) = \frac{\mathbb{E}[(Y(d_{i+1})-Y(d_i))\mathbf{1}_{\{D(1)\geq d_{i+1} > D(0)\}}]}{\mathbb{P}[D(1)\geq d_{i+1} > D(0)]}\, ,$$
é o efeito médio do tratamento de se manipular a causa $D$ de $d_i$ para $d_{i+1}$, dentre a subpopulação para a qual manipular hipoteticamente o valor do instrumento de $Z=0$ a $Z=1$ muda a dose do tratamento de $D(0) < d_{i+1}$ a $D(1) \geq d_{i+1}$. Para quais subpopulações a ponderação atribui maior peso ao efeito causal?
\vspace{0.5em}


\paragraph{Exercício 8 (\textit{Trade-off} quantidade-qualidade de filhos)} Seu objetivo é entender como a decisão de ter mais filhos impacta a alocação de recursos na educação do primeiro filho. A literatura de economia da família descreve um \textit{tradeoff} entre qualidade e quantidade na decisão de reprodução: ao escolherem ter mais filhos, famílias tendem a alocar menos recursos na educação dos filhos mais velhos. Seu objetivo é testar essa hipótese. Para isso, você se baseia em \citet{Ponczek2012} e usa dados familiares extraídos do censo brasileiro de 2010. Esses dados estão disponíveis em \texttt{censo\_data\_family\_size.csv}. Cada observação da base corresponde ao primeiro filho de uma família brasileira amostrada (famílias cujos primeiros filhos são gêmeos correspondem a mais de uma observação). Um dicionário com a descrição das variáveis da base está disponíbel em \texttt{dicionario.xlsx}. O conjunto de dados pode ser baixado \href{https://www.dropbox.com/scl/fi/2q5gv4c3uuxfehvzmjyjq/dados_censo.zip?rlkey=enbbd3iltf9gp7bic991zf4sw&st=zj6v1jg1&dl=0}{neste link}.

\begin{enumerate}[label=\roman*]
	\item Carregue os dados. Restrinja a base aos primeiros filhos de famílias com mais de um parto ($\text{\texttt{family\_number\_births}}\geq 2$). Qual é o tamanho familiar médio das famílias dos primeiros filhos? \textit{\textbf{Dica:} Instale o pacote \texttt{data.table} e use a função \texttt{fread} para carregar os dados mais rapidamente.}
	\item Para avaliar a relação prevista pela teoria econômica, você ajusta a seguinte especificação linear por MQO:
	
	\begin{equation*}
		\text{\texttt{first\_child\_years\_of\_education}}_i = \beta_0 + \beta_1 \cdot 	\text{\texttt{family\_number\_children}}_i+\epsilon_i
	\end{equation*}
	
	 Qual a interpretação da estimativa obtida? Ela é estatisticamente significante a 5\%? 
	
	\item Você acredita que a estimativa obtida no item anterior é crível para o efeito que se procura estimar? Por quê? \\ \\
	Como forma de mitigar as preocupações com a estimação obtida anteriormente, você se propõe a estimar o efeito de se ter mais filhos sobre a educação do primeiro filho instrumentando  $	\text{\texttt{family\_number\_children}}$ com $\text{\texttt{second\_birth\_ismultiplebirth}}$, uma variável binária que é igual a 1 se o segundo parto da família foi de mais de uma criança (gêmeos, trigêmeos etc.).
	
	\item Quantos segundos partos na base são partos de múltiplas crianças?
	
	\item Sob quais conjuntos de hipóteses a metodologia de variáveis instrumentais recupera um efeito causal de interesse, nos seguintes casos: (i) a resposta ao \textit{trade-off} educação e qualidade é homogênea na população; (2) a resposta ao \textit{trade-off} educação e qualidade é heterogênea na população. Quais são os efeitos causais identificados em cada um dos casos? Você consegue pensar em violações dessas hipóteses?
	
	\item Rode o primeiro estágio da metodologia de variáveis instrumentais. Qual a interepretação dos coeficientes obtidos? O instrumento é \textbf{relevante}, a 5\% de significância?
	
	\item Rode o segundo estágio. Qual a interpretação do coeficiente obtido? Ele é estatisticamente significante a 5\%? \textit{Observação:} você deve implementar erros padrão com cluster no nível da família (por quê?).
	
\item Devido ao padrão de amostragem, sabemos que cada indivíduo na amostra é representativo de \texttt{person\_weight} indivíduos na população. Reestime os modelos usando essa variável como peso. Os resultados mudam?
\end{enumerate}

\printbibliography
\end{document}